\documentclass[11pt, a4paper]{article}

% ── Codificación y Lenguaje ──────────────────────────────────────────────────
\usepackage[utf8]{inputenc}
\usepackage[T1]{fontenc}
\usepackage[spanish, es-tabla]{babel}

% ── Geometría ───────────────────────────────────────────────────────────────
\usepackage[top=2.5cm, bottom=2.5cm, left=2.5cm, right=2.5cm, headheight=24pt]{geometry}

% ── Matemáticas ─────────────────────────────────────────────────────────────
\usepackage{amsmath, amssymb, amsthm}
\usepackage{mathtools}

% ── Colores y Cajas ─────────────────────────────────────────────────────────
\usepackage{xcolor}
\usepackage[most]{tcolorbox}
\tcbuselibrary{skins, breakable, theorems}

% ── Tablas ──────────────────────────────────────────────────────────────────
\usepackage{booktabs}
\usepackage{array}
\usepackage{multirow}

% ── Listas ──────────────────────────────────────────────────────────────────
\usepackage{enumitem}

% ── Encabezado y Pie de Página ──────────────────────────────────────────────
\usepackage{fancyhdr}

% ── Secciones ───────────────────────────────────────────────────────────────
\usepackage{titlesec}

% ── Gráficos y Diagramas ────────────────────────────────────────────────────
\usepackage{tikz}
\usetikzlibrary{shapes.geometric, arrows.meta, positioning, calc}
\usepackage{graphicx}

% ── Columnas múltiples (formulario) ─────────────────────────────────────────
\usepackage{multicol}

% ── Espaciado ───────────────────────────────────────────────────────────────
\usepackage{setspace}
\setlength{\parskip}{4pt}
\setlength{\parindent}{0pt}

% ── Paleta GOP ──────────────────────────────────────────────────────────────
\definecolor{gopblue}{RGB}{31, 78, 121}
\definecolor{gopcyan}{RGB}{70, 130, 180}
\definecolor{gopgreen}{RGB}{0, 100, 0}
\definecolor{gopgreenlight}{RGB}{230, 255, 230}
\definecolor{goporange}{RGB}{200, 100, 0}
\definecolor{goporangelight}{RGB}{255, 245, 210}
\definecolor{gopred}{RGB}{160, 0, 0}
\definecolor{gopredlight}{RGB}{255, 230, 230}
\definecolor{gopgray}{RGB}{90, 90, 90}
\definecolor{gopgraylight}{RGB}{245, 245, 240}
\definecolor{gopbluedef}{RGB}{0, 70, 140}
\definecolor{gopbluedeflight}{RGB}{220, 235, 255}

% ── Hipervínculos y TOC ─────────────────────────────────────────────────────
\usepackage[colorlinks=true, linkcolor=gopblue, urlcolor=gopblue]{hyperref}

% ── Entornos tcolorbox (idénticos al resumen grande) ─────────────────────────
\newtcolorbox{definicion}[1]{
  enhanced, breakable,
  colback=gopbluedeflight, colframe=gopbluedef,
  fonttitle=\bfseries\small, title={Definicion: #1}, coltitle=white,
  attach boxed title to top left={yshift=-2mm, xshift=6pt},
  boxed title style={colback=gopbluedef, colframe=gopbluedef, sharp corners},
  sharp corners=south, top=4mm, before skip=6pt, after skip=6pt,
}
\newtcolorbox{formulas}[1]{
  enhanced, breakable,
  colback=gopgreenlight, colframe=gopgreen,
  fonttitle=\bfseries\small, title={Formulas: #1}, coltitle=white,
  attach boxed title to top left={yshift=-2mm, xshift=6pt},
  boxed title style={colback=gopgreen, colframe=gopgreen, sharp corners},
  sharp corners=south, top=4mm, before skip=6pt, after skip=6pt,
}
\newtcolorbox{tipprueba}[1]{
  enhanced, breakable,
  colback=goporangelight, colframe=goporange,
  fonttitle=\bfseries\small, title={TIP PRUEBA: #1}, coltitle=white,
  attach boxed title to top left={yshift=-2mm, xshift=6pt},
  boxed title style={colback=goporange, colframe=goporange, sharp corners},
  sharp corners=south, top=4mm, before skip=6pt, after skip=6pt,
}
\newtcolorbox{trampa}[1]{
  enhanced, breakable,
  colback=gopredlight, colframe=gopred,
  fonttitle=\bfseries\small, title={TRAMPA/ADVERTENCIA: #1}, coltitle=white,
  attach boxed title to top left={yshift=-2mm, xshift=6pt},
  boxed title style={colback=gopred, colframe=gopred, sharp corners},
  sharp corners=south, top=4mm, before skip=6pt, after skip=6pt,
}
\newtcolorbox{ejercicio}[1]{
  enhanced, breakable,
  colback=gopgraylight, colframe=gopgray,
  fonttitle=\bfseries\small\color{black}, title={Ejercicio: #1},
  attach boxed title to top left={yshift=-2mm, xshift=6pt},
  boxed title style={colback=gopgray!40, colframe=gopgray, sharp corners},
  sharp corners=south, top=4mm, before skip=6pt, after skip=6pt,
}
\newtcolorbox{resumentemario}{
  enhanced,
  colback=gopbluedeflight, colframe=gopbluedef,
  fonttitle=\bfseries, title={Estrategia de Estudio}, coltitle=white,
  attach boxed title to top left={yshift=-2mm, xshift=6pt},
  boxed title style={colback=gopbluedef, colframe=gopbluedef, sharp corners},
}
\newtcolorbox{estructuraprueba}{
  enhanced,
  colback=goporangelight, colframe=goporange,
  fonttitle=\bfseries, title={Estructura de la Prueba}, coltitle=white,
  attach boxed title to top left={yshift=-2mm, xshift=6pt},
  boxed title style={colback=goporange, colframe=goporange, sharp corners},
}

% ── Encabezado y pie ─────────────────────────────────────────────────────────
\pagestyle{fancy}
\fancyhf{}
\fancyhead[L]{\small\textbf{GOP ICS3213} --- Resumen \texttt{LITE} \texttt{I2} \texttt{2026}}
\fancyhead[C]{}
\fancyhead[R]{\small\nouppercase{\rightmark}}
\fancyfoot[L]{\footnotesize Compatible con \texttt{estudio\_i2.tex} --- Prioridad por ROI de certamen}
\fancyfoot[R]{\small\thepage}
\renewcommand{\headrulewidth}{0.4pt}
\renewcommand{\footrulewidth}{0.4pt}

% ── Estilo de secciones ──────────────────────────────────────────────────────
\titleformat{\section}
  {\color{gopblue}\Large\bfseries}{\color{gopblue}\thesection.}{0.5em}{}
  [\color{gopblue}\titlerule]
\titleformat{\subsection}
  {\color{gopblue}\normalsize\bfseries}{\color{gopblue}\thesubsection.}{0.5em}{}
\titleformat{\subsubsection}
  {\color{gopblue}\small\bfseries}{\color{gopblue}\thesubsubsection.}{0.5em}{}

\begin{document}

% ════════════════════════════════════════════════════════════════════════════
% PORTADA
% ════════════════════════════════════════════════════════════════════════════
\begin{titlepage}
  \centering
  \vspace*{1.5cm}
  {\Huge\bfseries\color{gopred} Resumen LITE \texttt{I2}\par}
  \vspace{0.3cm}
  {\LARGE\bfseries Gestión de Operaciones\par}
  \vspace{0.2cm}
  {\large ICS3213 --- PUC Chile --- \texttt{2026}\par}
  \vspace{0.8cm}
  \rule{\textwidth}{0.4pt}
  \vspace{0.4cm}

  {\normalsize \texttt{Miércoles 3 de Junio de 2026, 17:30 hrs (presencial)}\par}
  \vspace{0.2cm}
  {\small Prof.\ Alejandro Mac Cawley --- Rodrigo Carrasco\par}
  \vspace{0.2cm}
  {\small\textbf{Autor:} \href{https://cl.linkedin.com/in/fabian-ortega-llanten}{Fabián Ignacio Ortega Llantén}\par}
  \vspace{0.5cm}
  \rule{\textwidth}{0.4pt}
  \vspace{0.5cm}

  \begin{resumentemario}
    \textbf{Este es el resumen de emergencia.} Cubre solo lo que aparece en todos
    (o casi todos) los certámenes I2 (2020--2025), ordenado por \textbf{mayor puntaje
    y menor complejidad}. Si terminas este documento, puedes continuar con
    \texttt{estudio\_i2.tex} (los modulos tienen los mismos nombres).

    \vspace{4pt}
    \begin{tabular}{llccc}
      \toprule
      \textbf{Prioridad} & \textbf{Tema (Modulo)} & \textbf{Freq.} & \textbf{Pts} & \textbf{Dificultad} \\
      \midrule
      1 & PERT/CPM estocastico (Mod.\ 3)  & 5/5 & 15--20 & Baja \\
      2 & MRP + Silver-Meal/WW (Mod.\ 2)  & 5/5 & 15--20 & Baja--Media \\
      3 & MILP Planif.\ Agregada (Mod.\ 1) & 2/5 & 20     & Media \\
      4 & Colas G/G/1 -- Kingman (Mod.\ 4) & 1/5 & 20     & Media \\
      \midrule
      \multicolumn{2}{l}{\textit{Omitidos del lite (baja freq.\ o alta complejidad)}} & & & \\
      -- & Newsvendor / Inventarios         & 1/5 & 20 & Alta \\
      -- & Localizacion / Bodegas           & 1/5 & 20 & Alta \\
      \bottomrule
    \end{tabular}
  \end{resumentemario}

  \vspace{0.4cm}

  \begin{estructuraprueba}
    \begin{itemize}[leftmargin=1.2em, itemsep=3pt]
      \item \textbf{Etapa Digital (60 min, SIN apuntes):}
        V/F La Meta, seleccion multiple, 1 desarrollo corto.
      \item \textbf{Descanso de 5 minutos}
      \item \textbf{Etapa Desarrollo (60 min, APUNTES FISICOS + formulario):}
        Ejercicios de calculo --- \textbf{es la etapa que cubre este resumen LITE}.
        Solo papel. Escanear y subir a Canvas.
    \end{itemize}
  \end{estructuraprueba}

  \vfill
\end{titlepage}

\tableofcontents
\newpage

% ════════════════════════════════════════════════════════════════════════════
% MODULO 3: PERT/CPM (PRIORIDAD 1)
% ════════════════════════════════════════════════════════════════════════════
\section[Modulo 3: Proyectos PERT/CPM]{Modulo 3: Administracion de Proyectos (CPM / PERT)}
\markright{Modulo 3: PERT/CPM}

Aparece en \textbf{todos los certamenes I2 (2020--2025)}.
Procedimiento completamente mecanizable: aprende los 6 pasos y ganas 15--20 puntos.

% ── 3.1 CPM: Ruta Critica ─────────────────────────────────────────────────
\subsection{CPM: Calculo de Tiempos y Ruta Critica}

\begin{formulas}{Pasada Forward (ES/EF)}
  \begin{itemize}[leftmargin=1.2em, itemsep=2pt]
    \item $ES_i$ = Early Start = maximo de los $EF$ de todos los predecesores.
          (Para actividades sin predecesores: $ES = 0$.)
    \item $EF_i = ES_i + TE_i$
  \end{itemize}
\end{formulas}

\begin{formulas}{Pasada Backward (LS/LF)}
  \begin{itemize}[leftmargin=1.2em, itemsep=2pt]
    \item $LF_i$ = Late Finish = minimo de los $LS$ de todos los sucesores.
          (Para la ultima actividad: $LF = EF_{\max}$.)
    \item $LS_i = LF_i - TE_i$
    \item $\text{Holgura}_i = LS_i - ES_i = LF_i - EF_i$
  \end{itemize}
\end{formulas}

\begin{tipprueba}{Identificar la Ruta Critica en 3 segundos}
  La \textbf{ruta critica} es la cadena de actividades con holgura $= 0$.
  Su duracion es el $EF$ maximo del proyecto.
  Puede haber mas de una ruta critica.
\end{tipprueba}

% ── 3.2 PERT Estocastico ────────────────────────────────────────────────────
\subsection{PERT Estocastico: IC y Contratos}

\begin{formulas}{Estadisticas de la Ruta Critica}
  \[
    \mu_p = \sum_{i \in RC} TE_i
    \qquad
    \sigma_p^2 = \sum_{i \in RC} \sigma_i^2
    \qquad
    \sigma_p = \sqrt{\sigma_p^2}
  \]
  Donde $RC$ es el conjunto de actividades en la ruta critica.
\end{formulas}

\begin{formulas}{Intervalo de Confianza}
  \[
    IC_{1-\alpha} = \mu_p \;\pm\; Z_{\alpha/2} \cdot \sigma_p
  \]
  Valores usuales: $Z_{0.90} = 1.645$, $Z_{0.95} = 1.96$, $Z_{0.975} = 2.24$,
  $Z_{0.99} = 2.576$.
\end{formulas}

\begin{formulas}{Probabilidad de Terminar en Plazo $T$}
  \[
    Z = \frac{T - \mu_p}{\sigma_p}
    \qquad
    P(\text{terminar} \le T) = \Phi(Z)
    \qquad
    P(\text{terminar} > T) = 1 - \Phi(Z)
  \]
  Si $Z < 0$: el plazo $T$ es antes del tiempo esperado (probabilidad $< 50\%$).
\end{formulas}

\begin{formulas}{Evaluacion de Contratos (Valor Esperado)}
  \[
    VE_{\text{neto}} = \underbrace{B \cdot P(\text{terminar} \le T_B)}_{\text{VE bono}}
                    - \underbrace{P_{\text{pen}} \cdot P(\text{terminar} \ge T_P)}_{\text{VE penalizacion}}
  \]
  \textbf{Aceptar} si $VE_{\text{neto}} \ge 0$. \textbf{Rechazar} si $VE_{\text{neto}} < 0$.

  \medskip
  \textbf{Fecha de indiferencia} $T^*$: despejar $T$ tal que $VE_{\text{neto}} = 0$.
  \[
    B \cdot P(T \le T_B) = P_{\text{pen}} \cdot P(T \ge T^*)
    \;\Rightarrow\;
    P(T \ge T^*) = \frac{B \cdot \Phi(Z_B)}{P_{\text{pen}}}
    \;\Rightarrow\;
    T^* = \mu_p + Z^* \cdot \sigma_p
  \]
\end{formulas}

\begin{tipprueba}{Pasos para Resolver un Contrato}
  \begin{enumerate}[leftmargin=1.5em, itemsep=2pt]
    \item Calcular $\mu_p$ y $\sigma_p$ de la ruta critica.
    \item Para cada umbral ($T_B$, $T_P$): calcular $Z = (T - \mu_p)/\sigma_p$ y buscar $\Phi(Z)$.
    \item $VE_{\text{bono}} = B \cdot \Phi(Z_B)$; $VE_{\text{pen}} = P_{\text{pen}} \cdot (1 - \Phi(Z_P))$.
    \item Comparar: si $VE_{\text{bono}} > VE_{\text{pen}}$ aceptar.
    \item Para indiferencia: despejar $Z^*$ y luego $T^*$.
  \end{enumerate}
\end{tipprueba}

\begin{trampa}{Signos de Z: el error mas comun}
  $\Phi(Z)$ es la probabilidad de estar \textbf{por debajo} de $T$.
  Si el plazo es muy ajustado ($T < \mu_p$), entonces $Z < 0$ y $\Phi(Z) < 0.5$.
  Nunca pongas $\Phi(-1.36) = 0.9131$; es $1 - 0.9131 = 0.0869$.
\end{trampa}

% ── 3.3 Rutas Alternativas ─────────────────────────────────────────────────
\subsection[Rutas Alternativas]{Rutas Alternativas y Probabilidad de Volverse Critica}

\begin{tipprueba}{Procedimiento para Rutas Alternativas}
  \begin{enumerate}[leftmargin=1.5em, itemsep=2pt]
    \item Listar \textbf{todos} los caminos de inicio a fin del proyecto.
    \item Para cada ruta alternativa calcular su propio $\mu_r$ y $\sigma_r^2$
          (suma de $TE$ y $\sigma^2$ de sus actividades).
    \item Calcular $P(\text{ruta} > T_{\text{proyecto}})$ con
          $Z = (T_{\text{proyecto}} - \mu_r)/\sigma_r$.
    \item Una ruta alternativa ``se vuelve critica'' cuando supera la duracion
          esperada del proyecto.
  \end{enumerate}

  \textbf{Nota:} Las rutas pueden compartir actividades, por lo que sus duraciones
  no son independientes. La pauta solo pide la probabilidad individual de cada ruta.
\end{tipprueba}

% ── 3.4 Crashing ────────────────────────────────────────────────────────────
\subsection{Crashing: Reduccion de Duracion a Costo Minimo}

\begin{formulas}{Pendiente de Crashing}
  \[
    \text{Pendiente}_i = \frac{CC_i - CN_i}{TN_i - TC_i}
    \qquad [\$/\text{semana}]
  \]
  donde $CN$ = costo normal, $CC$ = costo crash, $TN$ = tiempo normal, $TC$ = tiempo crash.
  \textbf{Menor pendiente = mas barato acortar.}
\end{formulas}

\begin{tipprueba}{Algoritmo de Crashing con Valor Esperado (version certamen)}
  Cuando el crashing se evalua bajo incertidumbre (hay un bono por terminar antes):
  \begin{enumerate}[leftmargin=1.5em, itemsep=2pt]
    \item Para cada nivel de reduccion (0, 1, 2, \ldots semanas), identificar
          la combinacion mas barata de actividades a acortar en la ruta critica.
    \item Calcular el $VE_{\text{neto}} = \text{bono} \cdot P(\text{terminar antes}) - \text{costo crashing}$.
    \item Elegir el nivel de reduccion con mayor $VE_{\text{neto}}$ (puede ser 0 semanas).
    \item \textbf{Verificar rutas no criticas:} tras cada reduccion, calcular la duracion
          de las otras rutas. Si alguna supera a la ruta critica acortada, se vuelve critica
          y debes incluirla en el siguiente paso.
  \end{enumerate}
\end{tipprueba}

\begin{trampa}{No verificar rutas alternativas tras el crashing}
  Si reduces la ruta critica 2 semanas pero no checas las otras rutas,
  puedes terminar con una ruta no critica que ahora domina (se convierte en la nueva RC)
  y tu respuesta es incorrecta aunque el calculo del crashing este bien.
  \medskip
  \textbf{Regla rapida:} calcula la duracion de cada ruta no critica y compara con
  la nueva duracion de la RC. Si ninguna la supera, el analisis es valido.
\end{trampa}

\begin{trampa}{Crashing con Dos Rutas Criticas Simultaneas}
  Si hay dos rutas criticas al mismo tiempo, debes acortar ambas.
  Busca primero una actividad compartida (acortas las dos con un solo costo).
  Si no hay, debes acortar una actividad de cada ruta y sumar los costos.
\end{trampa}

% ════════════════════════════════════════════════════════════════════════════
% MODULO 2: MRP + LOTEO (PRIORIDAD 2)
% ════════════════════════════════════════════════════════════════════════════
\section[Modulo 2: MRP y Loteo]{Modulo 2: Planificacion de Corto Plazo y MRP}
\markright{Modulo 2: MRP y Loteo}

Aparece en \textbf{todos los certamenes I2 (2020--2025)}.
La tabla MRP es un procedimiento mecanico; Silver-Meal tiene un criterio de parada claro.

% ── 2.1 BOM y MRP ──────────────────────────────────────────────────────────
\subsection{Arbol BOM y Logica de la Tabla MRP}

\begin{definicion}{Lista de Materiales (BOM)}
  El BOM es un arbol jerarquico que indica cuantas unidades de cada componente
  se necesitan para fabricar una unidad del nivel superior. El \textbf{coeficiente}
  es el multiplicador. Si el producto final necesita 2 unidades de A, y A necesita
  3 de A1, entonces por unidad final se necesitan $2 \times 3 = 6$ unidades de A1.
\end{definicion}

\begin{formulas}{Las 6 filas de la tabla MRP}
  \begin{enumerate}[leftmargin=1.5em, itemsep=2pt]
    \item $GR_t$ (Requerimiento Bruto): demanda del periodo, o lanzamiento del nivel superior $\times$ coeficiente.
    \item $SR_t$ (Scheduled Receipts / Ordenes en Transito): ordenes ya colocadas con llegada en $t$.
    \item $I_t$ (Inventario al final de $t$): $I_t = I_{t-1} + SR_t + POR_t - GR_t$.
    \item $NR_t$ (Requerimiento Neto): $NR_t = \max(0,\; GR_t - SR_t - I_{t-1})$.
    \item $POR_t$ (Recepcion Planificada): cantidad que llega en $t$ para cubrir $NR_t$.
    \item $PORelease_t$ (Lanzamiento Planificado): $POR_t$ adelantado en el Lead Time ($LT$).
          $PORelease_{t-LT} = POR_t$.
  \end{enumerate}
\end{formulas}

\begin{tipprueba}{Procedimiento MRP Paso a Paso}
  \begin{enumerate}[leftmargin=1.5em, itemsep=2pt]
    \item Calcular $GR$ del nivel superior. Propagarlo hacia abajo multiplicando
          por el coeficiente BOM y usando el $PORelease$ como $GR$ del siguiente nivel.
    \item Para cada periodo: calcular $NR_t = \max(0,\; GR_t - SR_t - I_{t-1})$.
    \item Aplicar la politica de loteo ($POR_t$): L4L ($POR_t = NR_t$), Silver-Meal, WW.
    \item Actualizar inventario: $I_t = I_{t-1} + SR_t + POR_t - GR_t$.
    \item Adelantar $PORelease$ restando el Lead Time.
  \end{enumerate}
\end{tipprueba}

\begin{trampa}{Lead Time y Adelanto}
  Si $LT = 2$ semanas y la recepcion planificada es en la semana 6,
  el \textbf{lanzamiento} ($PORelease$) debe ir en la semana $6 - 2 = 4$.
  El error clasico es confundir la fila POR con la fila PORelease.
\end{trampa}

% ── 2.2 Silver-Meal ────────────────────────────────────────────────────────
\subsection{Heuristica de Silver-Meal}

\begin{formulas}{Costo Promedio por Periodo Silver-Meal}
  Para un lote lanzado en el periodo $t$ que cubre $k$ periodos:
  \[
    C(k) = \frac{S + H \sum_{j=1}^{k-1} j \cdot D_{t+j}}{k}
  \]
  donde $S$ = costo de setup, $H$ = costo de almacenamiento por unidad por periodo,
  $D_{t+j}$ = demanda en el periodo $t+j$.

  \medskip
  \textbf{Regla de parada:} calcular $C(1), C(2), C(3), \ldots$ hasta que
  $C(k+1) > C(k)$. El lote cubre los primeros $k$ periodos.
  Iniciar un nuevo lote a partir del periodo $t + k$.
\end{formulas}

\begin{trampa}{Periodos con demanda cero al inicio --- no empieces ahi}
  Si los primeros periodos tienen demanda cero, \textbf{saltate directamente al primer
  periodo con demanda positiva} y arranca Silver-Meal desde ahi.
  Empezar en un periodo con $D_t = 0$ genera $C(1) = S/1$ si haces setup (innecesario)
  o $C(1) = 0$ si no (comparacion degenerada). Ambos confunden sin aportar.
  \medskip
  La pauta de la Ay.\ 8 empieza en $t=0$ con demanda $=0$ y obtiene $C(1) = 0$,
  luego $C(2) = 73{,}5 > 0$ y ``para''. El resultado es correcto pero el razonamiento
  es innecesariamente confuso. La forma limpia: ignorar $t=1$ y arrancar en $t=2$.
\end{trampa}

\begin{ejercicio}{Silver-Meal: esquema rapido}
  Demanda: $[30, 120, 180, 150]$. $S = 120$, $H = 0.9$.

  \medskip
  \textbf{Lote 1 (desde periodo 1):}
  \begin{itemize}[leftmargin=1.5em, itemsep=1pt]
    \item $C(1) = 120 / 1 = 120$
    \item $C(2) = (120 + 1 \cdot 0.9 \cdot 120) / 2 = 228 / 2 = 114$ $\to$ baja, seguir.
    \item $C(3) = (228 + 2 \cdot 0.9 \cdot 180) / 3 = 552 / 3 = 184$ $\to$ sube. \textbf{Parar en $k=2$.}
  \end{itemize}
  Lote 1: cubre periodos 1 y 2 $\Rightarrow$ $30 + 120 = 150$ unidades.

  \medskip
  \textbf{Lote 2 (desde periodo 3):}
  \begin{itemize}[leftmargin=1.5em, itemsep=1pt]
    \item $C(1) = 120$
    \item $C(2) = (120 + 1 \cdot 0.9 \cdot 150) / 2 = 255 / 2 = 127.5$ $\to$ sube. \textbf{Parar en $k=1$.}
  \end{itemize}
  Lote 2: solo periodo 3 $\Rightarrow$ 180 unidades.

  \textbf{Lote 3}: periodo 4 $\Rightarrow$ 150 unidades.
\end{ejercicio}

% ── 2.3 Wagner-Whitin ──────────────────────────────────────────────────────
\subsection{Algoritmo de Wagner-Whitin (WW)}

\begin{formulas}{Programacion Dinamica WW}
  Sea $f(t)$ el costo minimo para cubrir la demanda de los periodos $1, \ldots, t$:
  \[
    f(t) = \min_{1 \le j \le t} \left[ S + c(j, t) + f(j-1) \right]
  \]
  donde $c(j, t)$ es el costo de almacenamiento si el lote lanzado en $j$ cubre
  hasta $t$:
  \[
    c(j, t) = H \sum_{i=j+1}^{t} (i - j) \cdot D_i
  \]
  Condicion inicial: $f(0) = 0$.
  Se reconstruye la solucion trazando hacia atras los $j^*$ que minimizan cada $f(t)$.
\end{formulas}

\begin{tipprueba}{Tabla WW Paso a Paso}
  \begin{enumerate}[leftmargin=1.5em, itemsep=2pt]
    \item Construir la tabla de costos $c(j,t)$ para todo $j \le t$.
    \item Para $t = 1$: $f(1) = S + 0 = S$ (el lote cubre solo el periodo 1).
    \item Para cada $t$ siguiente, evaluar todas las opciones $j$ y tomar la minima.
    \item Al final, reconstruir los lotes trazando $j^*(t) \to j^*(j^*-1) \to \ldots$
  \end{enumerate}
\end{tipprueba}

% ════════════════════════════════════════════════════════════════════════════
% MODULO 1: MILP PLANIFICACION AGREGADA (PRIORIDAD 3)
% ════════════════════════════════════════════════════════════════════════════
\section[Modulo 1: Planificacion Agregada MILP]{Modulo 1: Planificacion Agregada --- Modelo MILP}
\markright{Modulo 1: Planificacion Agregada}

Aparecio en los \textbf{dos certamenes mas recientes (2024 y 2025)}.
La pregunta siempre tiene la misma estructura: (i) variables + FO + restricciones base;
(ii) agregar condiciones complejas con binarias. Si memorizas la \textbf{receta de 5 preguntas}
puedes plantear cualquier modelo sin haberlo visto antes.

% ── 1.0 Receta general ────────────────────────────────────────────────────
\subsection{La Receta: 5 Preguntas para Plantear Cualquier MILP}

\begin{tipprueba}{Receta MILP --- Hazte estas 5 preguntas en orden}
  \begin{enumerate}[leftmargin=1.5em, itemsep=4pt]

    \item \textbf{?`Que DECIDO cada periodo?} $\to$ Variables de decision.
      \begin{itemize}[leftmargin=1.2em, itemsep=1pt]
        \item Continuas ($\ge 0$): cuanto produzco ($P$), cuanto guardo ($I$),
              cuanto despacho ($S$), cuantas horas extra ($HE$), cuantos trabajadores ($TA$).
        \item Binarias ($\in \{0,1\}$): ?`opera la planta? ($X$), ?`arranca esta semana? ($Y$).
      \end{itemize}
      \textit{Regla practica:} necesitas una variable por cada cantidad que el modelo
      debe elegir. Si el enunciado dice ``decide si'' $\to$ binaria; si dice ``decide cuanto'' $\to$ continua.

    \item \textbf{?`Que CUESTA?} $\to$ Funcion objetivo $= \min \sum (\text{costo unitario} \times \text{variable})$.
      Recorre el enunciado buscando: produccion, inventario, despacho/envio, horas extra,
      salarios, contratacion/despido, activacion de planta, ventas perdidas.
      Cada uno es un termino en la FO.

    \item \textbf{?`Que DEBO satisfacer?} $\to$ Restriccion de demanda (una por cliente y periodo).
      \[
        \sum_n S_{n,i,t} + NS_{i,t} = D_{it} \quad \forall i,t
        \qquad \text{(si hay ventas perdidas)}
      \]
      \[
        \sum_n S_{n,i,t} \ge D_{it} \quad \forall i,t
        \qquad \text{(si no hay ventas perdidas)}
      \]

    \item \textbf{?`Que PUEDO hacer?} $\to$ Restricciones de capacidad.
      La produccion esta acotada por los recursos disponibles.
      Si hay una variable binaria $X_{n,t}$, la capacidad se ``enciende'' con ella:
      \[
        P_{n,t} \le \text{capacidad}_n \cdot X_{n,t} \quad \forall n,t
      \]

    \item \textbf{?`Como EVOLUCIONA el inventario?} $\to$ Restriccion de balance (una por planta y periodo).
      \[
        I_{n,t} = I_{n,t-1} + P_{n,t} - \sum_i S_{n,i,t} \quad \forall n,t
      \]
      Incluir condiciones de borde: $I_{n,0} = I_0^{\text{dato}}$, $I_{n,T} = I_T^{\text{requerido}}$.

  \end{enumerate}

  \medskip
  \textbf{Parte (ii) del certamen: condiciones complejas.}
  Lee cada condicion en lenguaje natural e identifica el patron:
  \begin{itemize}[leftmargin=1.2em, itemsep=2pt]
    \item ``Si opera, entonces al menos $L$ unidades'' $\to P_{n,t} \ge L \cdot X_{n,t}$
    \item ``Solo produce si opera'' $\to P_{n,t} \le M \cdot X_{n,t}$ (Big-M)
    \item ``Si arranca, paga costo fijo'' $\to$ necesitas $Y_{n,t}$ con $Y_{n,t} \ge X_{n,t} - X_{n,t-1}$
    \item ``A lo mas el $\alpha$\% insatisfecho'' $\to \sum_t NS_{i,t} \le \alpha \cdot \sum_t D_{it}$
    \item ``Si A ocurre, B debe ocurrir'' $\to X^A \le X^B$ (o $X^A + X^B \le 1$ si son excluyentes)
  \end{itemize}
\end{tipprueba}

% ── 1.1 Variables tipicas ──────────────────────────────────────────────────
\subsection{Variables de Decision Tipicas}

\begin{definicion}{Catalogo de Variables (Mod.\ 1)}
  \begin{itemize}[leftmargin=1.2em, itemsep=2pt]
    \item $P_{n,t} \ge 0$: produccion en planta $n$, semana $t$.
    \item $I_{n,t} \ge 0$: inventario en planta $n$ al final de $t$.
    \item $S_{n,i,t} \ge 0$ (o $DE$): despacho de planta $n$ a cliente $i$ en $t$.
    \item $NS_{i,t} \ge 0$: demanda insatisfecha del cliente $i$ en $t$ (si se permite).
    \item $TA_{n,t} \ge 0$: trabajadores en planta $n$ semana $t$.
    \item $TCA_{n,t}, TDA_{n,t} \ge 0$: contratados y despedidos al inicio de $t$.
    \item $HE_{n,t} \ge 0$: horas extra en planta $n$, semana $t$.
    \item $X_{n,t} \in \{0,1\}$: 1 si la planta $n$ opera en $t$.
    \item $Y_{n,t} \in \{0,1\}$: 1 si la planta $n$ \textbf{se activa} en $t$
          (estaba inactiva en $t-1$).
  \end{itemize}
\end{definicion}

% ── 1.2 FO tipica ──────────────────────────────────────────────────────────
\subsection{Funcion Objetivo Tipica}

\begin{formulas}{Componentes del Costo Total (Minimizar)}
  \[
    \min \; \sum_t \sum_n \Big(
      CP_n P_{n,t} + H_n I_{n,t} + A_n Y_{n,t} + CE_n HE_{n,t}
    \Big)
  \]
  \[
    + \sum_t \sum_n \Big(
      CT\, TA_{n,t} + CC\, TCA_{n,t} + CDS\, TDA_{n,t}
      + \textstyle\sum_i C_{in} S_{n,i,t}
    \Big)
    + \sum_t \sum_i P_i \cdot NS_{i,t}
  \]
  No todas las variables aparecen en todos los modelos; incluir solo las del enunciado.
\end{formulas}

% ── 1.3 Restricciones clave ────────────────────────────────────────────────
\subsection{Restricciones Clave}

\begin{formulas}{Restricciones Fundamentales}
  \textbf{Balance de inventario:}
  \[
    I_{n,t} = I_{n,t-1} + P_{n,t} - \sum_i S_{n,i,t} \quad \forall n, t
  \]

  \textbf{Satisfaccion de demanda} (con o sin ventas perdidas):
  \[
    \sum_n S_{n,i,t} + NS_{i,t} = D_{it} \quad \forall i, t
    \qquad \text{(con ventas perdidas)}
  \]
  \[
    \sum_n S_{n,i,t} \ge D_{it} \quad \forall i, t
    \qquad \text{(sin ventas perdidas)}
  \]

  \textbf{Capacidad de produccion} (horas normales + extras):
  \[
    P_{n,t} \le PR_n \cdot (HN_n \cdot X_{n,t} + HE_{n,t}) \quad \forall n, t
  \]
  \[
    HE_{n,t} \le HE_{\max,n} \cdot X_{n,t} \quad \forall n, t
  \]

  \textbf{Dinamica de fuerza laboral:}
  \[
    TA_{n,t} = TA_{n,t-1} + TCA_{n,t} - TDA_{n,t} \quad \forall n, t
  \]

  \textbf{Activacion de planta} (detectar arranque):
  \[
    Y_{n,t} \ge X_{n,t} - X_{n,t-1} \quad \forall n, t
  \]
\end{formulas}

\begin{tipprueba}{Condiciones Complejas con Big-M}
  Estas aparecen frecuentemente en la parte (ii) de la pregunta:
  \begin{itemize}[leftmargin=1.2em, itemsep=3pt]
    \item \textbf{Produccion minima si opera:}
      $P_{n,t} \ge L_n \cdot X_{n,t}$
    \item \textbf{Turno completo obligatorio:}
      $P_{n,t} \ge PR_n \cdot 40 \cdot X_{n,t}$
    \item \textbf{Produccion 0 si no opera (Big-M):}
      $P_{n,t} \le M \cdot X_{n,t}$
    \item \textbf{Limite de demanda insatisfecha:}
      $\displaystyle\sum_t NS_{i,t} \le A_i \cdot \sum_t D_{it}$
    \item \textbf{Relacion entre variables binarias:} si $X=1$ implica $Z=1$,
      usar $X_{n,t} \le Z_{n,t}$ o $X_{n,t} + Z_{n,t} \le 1$ segun la logica.
  \end{itemize}
\end{tipprueba}

\begin{trampa}{Plan Chase: el supuesto de redondeo que las pautas no declaran}
  En la practica, los panaderos son enteros. Al calcular ``trabajadores necesarios $= D_t / \text{cap}$''
  suele salir un numero fraccionario. Las pautas redondean ($7{,}5 \to 8$), pero eso
  significa que se \textbf{produce mas de lo demandado} (80 unidades extra que la pauta ignora).
  Si el enunciado dice ``producir exactamente la demanda'', usa el numero exacto aunque sea fraccionario.
  Si dice ``trabajadores enteros'', redondea y ajusta la produccion en consecuencia.
  \medskip
  \textbf{En la prueba:} si hay ambiguedad, indica el supuesto que estas asumiendo.
  La pauta te da el beneficio de la duda si el calculo es consistente con tu supuesto.
\end{trampa}

\begin{trampa}{El costo de activacion va en $Y_{n,t}$, no en $X_{n,t}$}
  $A_n \cdot X_{n,t}$ cobra el costo cada semana que opera.
  $A_n \cdot Y_{n,t}$ solo cobra cuando la planta \textbf{arranca} (transicion $0 \to 1$).
  Verifica siempre que la restriccion $Y_{n,t} \ge X_{n,t} - X_{n,t-1}$ este en el modelo.
\end{trampa}

% ════════════════════════════════════════════════════════════════════════════
% MODULO 4: COLAS G/G/1 (PRIORIDAD 4)
% ════════════════════════════════════════════════════════════════════════════
\section[Modulo 4: Variabilidad y Colas G/G/1]{Modulo 4: Variabilidad en las Operaciones (Colas G/G/1)}
\markright{Modulo 4: Variabilidad -- Colas G/G/1}

Aparecio en 1/5 certamenes (2022). Alta probabilidad de aparecer este ano dado que
el modulo es parte del temario oficial. Las formulas son mecanicas.

% ── 4.1 Formulas G/G/1 ─────────────────────────────────────────────────────
\subsection{Formula de Kingman (VUT)}

\begin{formulas}{Formula de Kingman para Tiempo en Cola G/G/1}
  \[
    CT_q = \underbrace{\frac{C_a^2 + C_s^2}{2}}_{\text{Variabilidad}}
           \cdot \underbrace{\frac{\rho}{1 - \rho}}_{\text{Utilizacion}}
           \cdot \underbrace{ST}_{\text{Tiempo de Servicio}}
  \]
  donde:
  \begin{itemize}[leftmargin=1.5em, itemsep=2pt]
    \item $\rho = \lambda / \mu$ (utilizacion, debe ser $< 1$ para estabilidad).
    \item $C_a$ = coeficiente de variacion de llegadas $ = \sigma_a / (1/\lambda)$.
    \item $C_s$ = coeficiente de variacion del servicio $ = \sigma_s \cdot \mu$.
    \item $ST = 1/\mu$ = tiempo promedio de servicio.
  \end{itemize}

  \medskip
  \textbf{Caso especial M/M/1:} $C_a = C_s = 1$ $\Rightarrow$
  $CT_q = \frac{1}{\mu - \lambda} = \frac{\rho}{\mu(1-\rho)}$.
\end{formulas}

\begin{formulas}{Metricas de la Cola}
  \[
    W_q = CT_q \quad [\text{tiempo en cola}]
    \qquad
    W = W_q + \frac{1}{\mu} \quad [\text{tiempo total en sistema}]
  \]
  \[
    L_q = \lambda \cdot W_q \quad [\text{clientes en cola}]
    \qquad
    L = \lambda \cdot W \quad [\text{clientes en sistema}]
  \]
  (Ley de Little: $L = \lambda \cdot W$)
\end{formulas}

% ── 4.2 Optimizacion de capacidad ──────────────────────────────────────────
\subsection{Optimizacion de Capacidad}

\begin{tipprueba}{Minimizar costo total: capacidad + espera}
  El costo total suele tener la forma:
  \[
    CT(\mu) = c_\mu \cdot \mu + c_w \cdot \lambda \cdot W_q(\mu)
  \]
  \begin{enumerate}[leftmargin=1.5em, itemsep=2pt]
    \item Expresar $W_q$ en funcion de $\mu$ usando Kingman.
    \item Derivar $CT(\mu)$ respecto a $\mu$ e igualar a cero.
    \item Verificar que $\rho = \lambda/\mu < 1$ en la solucion.
  \end{enumerate}
  Si la derivacion es compleja, probar con valores enteros y comparar costos.
\end{tipprueba}

% ── 4.3 Propagacion de variabilidad en serie ──────────────────────────────
\subsection{Propagacion de Variabilidad en Serie}

\begin{formulas}{Coeficiente de Variacion de Salida (Aproximacion)}
  \[
    C_{d}^2 \approx (1 - \rho^2) \cdot C_a^2 + \rho^2 \cdot C_s^2
  \]
  $C_d$ es el coeficiente de variacion de las salidas de la estacion,
  que se convierte en el $C_a$ de la siguiente estacion en serie.

  \medskip
  \textbf{Interpretacion:}
  \begin{itemize}[leftmargin=1.2em, itemsep=1pt]
    \item Alta utilizacion ($\rho \to 1$): $C_d \approx C_s$ (el servicio domina la variabilidad de salida).
    \item Baja utilizacion ($\rho \to 0$): $C_d \approx C_a$ (las llegadas dominan).
  \end{itemize}
\end{formulas}

\begin{trampa}{Cuando M/M/1 no aplica}
  Si el enunciado da $C_a \ne 1$ o $C_s \ne 1$, debes usar Kingman completo.
  M/M/1 es solo el caso $C_a = C_s = 1$ (llegadas y servicio exponenciales).
\end{trampa}

% ════════════════════════════════════════════════════════════════════════════
% FORMULARIO RAPIDO
% ════════════════════════════════════════════════════════════════════════════
\newpage
\section*{Formulario Rapido --- Resumen LITE}
\markright{Formulario Rapido}
\addcontentsline{toc}{section}{Formulario Rapido}

\begin{multicols}{2}

\subsubsection*{PERT / CPM}
\begin{itemize}[leftmargin=1.2em, itemsep=1pt, label={$\bullet$}]
  \item $ES_i = \max(EF_{\text{predecesores}})$
  \item $EF_i = ES_i + TE_i$
  \item $LS_i = LF_i - TE_i$
  \item $LF_i = \min(LS_{\text{sucesores}})$
  \item Holgura $= LS - ES$; RC: holgura $= 0$
  \item $\mu_p = \sum_{RC} TE_i$; $\sigma_p^2 = \sum_{RC} \sigma_i^2$
  \item $Z = (T - \mu_p)/\sigma_p$; $P(\le T) = \Phi(Z)$
  \item $IC_{95\%} = \mu_p \pm 1.96\,\sigma_p$
  \item $VE = B \cdot \Phi(Z_B) - P\cdot(1-\Phi(Z_P))$
  \item Pendiente crashing $= (CC-CN)/(TN-TC)$
\end{itemize}

\subsubsection*{MRP}
\begin{itemize}[leftmargin=1.2em, itemsep=1pt, label={$\bullet$}]
  \item $NR_t = \max(0,\; GR_t - SR_t - I_{t-1})$
  \item $I_t = I_{t-1} + SR_t + POR_t - GR_t$
  \item L4L: $POR_t = NR_t$
  \item $PORelease_{t-LT} = POR_t$
\end{itemize}

\subsubsection*{Silver-Meal}
\begin{itemize}[leftmargin=1.2em, itemsep=1pt, label={$\bullet$}]
  \item $C(k) = \left[S + H\sum_{j=1}^{k-1} j\,D_{t+j}\right]/k$
  \item Parar cuando $C(k+1) > C(k)$
\end{itemize}

\subsubsection*{Wagner-Whitin}
\begin{itemize}[leftmargin=1.2em, itemsep=1pt, label={$\bullet$}]
  \item $f(t) = \min_{j \le t}[S + c(j,t) + f(j-1)]$
  \item $c(j,t) = H\sum_{i=j+1}^{t}(i-j)D_i$; $f(0)=0$
\end{itemize}

\subsubsection*{Colas G/G/1 (Kingman)}
\begin{itemize}[leftmargin=1.2em, itemsep=1pt, label={$\bullet$}]
  \item $\rho = \lambda/\mu$ (debe ser $<1$)
  \item $CT_q = \frac{C_a^2+C_s^2}{2}\cdot\frac{\rho}{1-\rho}\cdot\frac{1}{\mu}$
  \item $W = CT_q + 1/\mu$; $L = \lambda\,W$
  \item M/M/1: $CT_q = \frac{1}{\mu-\lambda}$
  \item $C_d^2 \approx (1-\rho^2)C_a^2 + \rho^2 C_s^2$
\end{itemize}

\subsubsection*{MILP: Restricciones Clave}
\begin{itemize}[leftmargin=1.2em, itemsep=1pt, label={$\bullet$}]
  \item Balance: $I_{n,t} = I_{n,t-1} + P_{n,t} - \sum_i S_{n,i,t}$
  \item Demanda: $\sum_n S_{n,i,t} + NS_{i,t} = D_{it}$
  \item Capacidad: $P_{n,t} \le PR_n(HN_n X_{n,t} + HE_{n,t})$
  \item Activacion: $Y_{n,t} \ge X_{n,t} - X_{n,t-1}$
  \item Min prod.: $P_{n,t} \ge L_n X_{n,t}$
  \item Big-M: $P_{n,t} \le M\,X_{n,t}$
\end{itemize}

\subsubsection*{Tabla Normal Estandar (valores clave)}
\begin{tabular}{cc|cc}
  \toprule
  $Z$ & $\Phi(Z)$ & $Z$ & $\Phi(Z)$ \\
  \midrule
  0.00 & 0.5000 & 1.28 & 0.8997 \\
  0.43 & 0.6664 & 1.36 & 0.9131 \\
  0.52 & 0.6985 & 1.45 & 0.9265 \\
  0.64 & 0.7389 & 1.53 & 0.9370 \\
  0.68 & 0.7517 & 1.64 & 0.9495 \\
  0.84 & 0.7995 & 1.96 & 0.9750 \\
  1.00 & 0.8413 & 2.04 & 0.9793 \\
  1.02 & 0.8461 & 2.33 & 0.9901 \\
  \bottomrule
\end{tabular}

\end{multicols}

\end{document}
